\documentclass[12pt]{article}
\usepackage[spanish]{babel}
\usepackage[utf8]{inputenc}
\usepackage[T1]{fontenc}
\usepackage{amsmath, amssymb}
\usepackage{geometry}
\usepackage{enumitem}
\usepackage{needspace}
\usepackage{tikz}
\usetikzlibrary{babel}

\geometry{letterpaper, margin=2cm}
\setlength{\parindent}{0pt}
\setlength{\parskip}{6pt}

\newcommand{\puntografico}{0.09}

\begin{document}

\begin{center}
  {\Large \textbf{Borrador de items}}\\
  {\large Bloque 1. Geometria - Afirmacion 6}\\[4pt]
  Prueba Nacional Estandarizada de Matematica, Secundaria 2026
\end{center}

\section*{Afirmacion}
Resuelve problemas relacionados con simetrias en diversas figuras.

\section*{Items propuestos}

\begin{enumerate}[label=\textbf{\arabic*.}, itemsep=1.05cm]

\Needspace{0.82\textheight}
\item
En una institucion educativa se diseno una pieza decorativa para colocarla en la entrada del gimnasio. La siguiente figura representa la forma de esa pieza y algunas rectas trazadas sobre ella:

\begin{center}
\begin{tikzpicture}[scale=0.72]
  \path[use as bounding box] (-4.2,-3.1) rectangle (4.3,4.1);
  \fill[gray!18] (0,3) -- (2.5,0.8) -- (1.6,-2) -- (0,-1.1) -- (-1.6,-2) -- (-2.5,0.8) -- cycle;
  \draw[very thick] (0,3) -- (2.5,0.8) -- (1.6,-2) -- (0,-1.1) -- (-1.6,-2) -- (-2.5,0.8) -- cycle;
  \draw[thick, dashed] (0,-2.7) -- (0,3.6);
  \draw[thick, dashed] (-3.5,0.35) -- (3.5,0.35);
  \draw[thick, dashed] (-2.8,-2.4) -- (2.8,3.2);
  \draw[thick, dashed] (-2.8,3.2) -- (2.8,-2.4);
  \node[above] at (0,3.6) {$h$};
  \node[right] at (3.5,0.35) {$j$};
  \node[above right] at (2.8,3.2) {$k$};
  \node[below right] at (2.8,-2.4) {$n$};
\end{tikzpicture}
\end{center}

De acuerdo con la informacion anterior, \textquestiondown cual de esas rectas corresponde a un eje de simetria de la pieza decorativa?

\begin{enumerate}[label=\Alph*)]
  \item $h$
  \item $j$
  \item $k$
  \item $n$
\end{enumerate}

\Needspace{0.88\textheight}
\item
En el diseno de un parque se desea ubicar dos postes de iluminacion de manera simetrica con respecto a un sendero rectilineo. En el plano, el sendero esta representado por la recta $r$, y el poste $P$ esta ubicado como se muestra en la siguiente representacion grafica, cuyas unidades estan en metros:

\begin{center}
\begin{tikzpicture}[scale=0.48]
  \path[use as bounding box] (-1.2,-1.1) rectangle (12.4,8.7);
  \draw[<->, thick] (-0.5,0) -- (11.8,0);
  \draw[<->, thick] (0,-0.5) -- (0,8.3);
  \node[anchor=west] at (12.05,-0.2) {$x$};
  \node[anchor=south east] at (0.85,8.45) {$y$};
  \draw[very thick] (6,-0.2) -- (6,8.1);
  \node[anchor=south west] at (6,8.1) {$r$};
  \draw[dashed] (2,0) -- (2,3);
  \draw[dashed] (0,3) -- (2,3);
  \draw[dashed] (6,0) -- (6,3);
  \fill (2,3) circle (\puntografico) node[above left] {$P$};
  \node[anchor=north] at (2,0) {$2$};
  \node[anchor=north] at (6,0) {$6$};
  \node[left] at (0,3) {$3$};
\end{tikzpicture}
\end{center}

Si el poste $Q$ debe ubicarse en el punto simetrico de $P$ con respecto a la recta $r$, entonces \textquestiondown cuales son las coordenadas de $Q$?

\begin{enumerate}[label=\Alph*)]
  \item $(4,3)$
  \item $(6,3)$
  \item $(8,3)$
  \item $(10,3)$
\end{enumerate}

\end{enumerate}

\section*{Clave y justificacion breve}

\begin{enumerate}[label=\textbf{\arabic*.}]
  \item \textbf{A.} La recta $h$ divide la figura en dos partes congruentes que coinciden al reflejarse una sobre la otra; por eso corresponde a un eje de simetria.
  \item \textbf{D.} La recta $r$ corresponde a $x=6$. El punto $P(2,3)$ esta a $4$ unidades a la izquierda de esa recta, por lo que su simetrico debe estar $4$ unidades a la derecha: $Q(10,3)$.
\end{enumerate}

\end{document}
